\documentclass{comjnl}

\usepackage{amsmath}
\usepackage{siunitx}
\usepackage{enumitem}

%\copyrightyear{2009} \vol{00} \issue{0} \DOI{000}

\begin{document}


\title{In Vitro Anti-Inflammatory Activity Assessment Using Albumin Denaturation Inhibition Assay}
\author{Sayed Huzaifa Mumit ^{1,2,*} \\ 
\textit{1. The Bangladesh Bioscience Research Group (BBRG)} \\
\textit{2. University of Science and Technology Chittagong (USTC)} }
\email{huzaifa.dbb@ustc.ac.bd \\ huzaifamumit2004@gmail.com}

\shortauthors{S.H. Mumit}




%\category{C.2}{Computer Communication Networks}{Computer Networks}
%\category{C.4}{Performance of Systems}{Analytical Models}
%\category{G.3}{Stochastic Processes}{Queueing Systems}
%\terms{Internet Technologies, E-Commerce}
\keywords{Anti-inflammatory activity; Protein denaturation inhibition; 
Bovine serum albumin; In vitro assay; Spectrophotometry analysis; Phytochemical extracts}


\begin{abstract}
Protein denaturation is a well-recognized process associated with inflammatory responses, in which structural alterations of proteins may lead to the formation of auto-antigens and subsequent immune activation. The inhibition of protein denaturation is therefore widely used as an in vitro indicator for evaluating the anti-inflammatory potential of bioactive compounds. 

The present protocol describes a spectrophotometric method for assessing anti-inflammatory activity based on the inhibition of heat-induced denaturation of bovine serum albumin (BSA). In this assay, the protein solution is incubated with test samples, such as plant extracts or synthetic compounds, under controlled physiological conditions in a phosphate buffer system. The reaction mixture is subsequently subjected to thermal denaturation, and the resulting turbidity caused by protein aggregation is measured at 660\,nm using a spectrophotometer. 

Compounds exhibiting anti-inflammatory activity stabilize the protein structure and reduce denaturation, resulting in lower turbidity and absorbance compared with the untreated control. The percentage inhibition of protein denaturation is calculated relative to the control, providing an estimate of anti-inflammatory activity. This protocol offers a rapid, reproducible, and cost-effective approach for the preliminary screening of natural products and synthetic compounds prior to in vivo investigation.

\end{abstract}

\maketitle


\section{Introduction}
The protocol always to determine the phytochemical activity against the disease or specific inflammations in the human body. 
Here according to the universally accepted protocol will be discussed before the project proceed in the the laboratory.
There is some of the stoichiometry is require for the measure concentration and assay requirement.

\section{Background}
During Inflammation Human body and tissue experiences several conditions that increases the body or specific cell/ tissue temperature where physiological says 37 \unit{\degree} C. Reactive Oxidative Stress happens into the cellular function which produces free radical of Oxygen during ATP generation process and carrier protein dysfunction and GSH insufficiency. pH changes to  the acidic nature and Reactive metabolites activate.
These conditions allow the cell to become auto-antigenic and trigger to the immune signaling and stimulate synthesis of some cells they are immune cell : IL-1\unit{\beta}, pro-inflammatory TNF-\unit{\alpha}, cytokine IL-6 prostaglandin via Cox Pathway.
Usages of Bovine Serum Albumin (BSA) in In Vitro experiment because of cost effective and structural well characterized like human body cell protein, highly stable but heat sensitive, functionally similar to the human body. Albumin generally blood protein. 

Heat Induced Denaturation technique mimics functional condition in vitro.
Physiological equilibrium at 37 \unit{\degree} C.  
Controlled Denaturation at 57 \unit{\degree} C. at 57 \unit{\degree} C hydrogen bond breaks, hydrophobic region exposes,  aggregation begins, turbid image appears for spectrophotometer.

When the sample of the extract added to the protein solution there a anti-inflammatory activity to be observed by the bio active compounds within the sample which basically stabilize the protein conformation, prevent unfolding, reduce the aggregation and lower absorbance at 660 nm in spectrophotometer. The principle that higher turbid = higher absorbance which determine protein inflammation stage and lower turbid = lower absorbance which determine anti-inflammation activity.

The calculation for the measurement of the anti-inflammation activity how much with the equation
\begin{center}
    \%Inhibition = $((A_c - A_s)/ A_c ) *100$ \\
  A$_{c}$ = Fully denatured control \\
  A$_{s}$ = Protected Sample
\end{center}


\section{Materials and Methods}  \label{Performance}

\subsection{Materials}
 To determine the activity there \textbf{plant extracts} are required of the specific solvent.
 In Vitro process is artificial environment to create as like the human body within. 
Thus, some of the special reagents are required whether some of them may will require to prepare in the lab by mixing of raw chemicals. Anti-inflammatory assay In Vitro assay requires the reagent 
 \\
 \\
\textbf{Bovine Serum Albumin (BSA), Phosphate Buffer Solution(PBS),}.
All kind of experiments require triplicate to get best results.
\textbf{Control, Test, Standard} these will be examined triple times. \\
\textbf{Major}
\begin{itemize}
    \item [1] Concentration requires : 7  (triplicate : 21 times) 
    \item [2] BSA  (Control/Negative) 35mg
    \item [3] PBS  (Ph 6.3-6.5) 5ml
    \item [4] Plant Extract 50 mg
\end{itemize}


\textbf{Equipments}
\begin{itemize}
    \item Test tubes 42 
    \item Solvent 200ml
    \item Cuvette 2 pieces ( 3.5ml)
    \item 500 \unit{\micro\liter} Pipette 
    \item Gloves, Tissue, Mask
    \item Distilled Water 5L
\end{itemize} \\ 

\subsection{Method}
\subsubsection{Concentration Preparation Using Serial Dilution}
\qty{1000}{\micro\gram\per\milli\litre} will be prepared. Sample and standard both concentration requires 1 mg/1ml. 
Concentrations are
\begin{enumerate}
    \item 1000 \unit{\micro\gram\per\milli\litre}
    \item 500 \unit{\micro\gram\per\milli\litre}
    \item 250 \unit{\micro\gram\per\milli\litre}
    \item 125 \unit{\micro\gram\per\milli\litre}
    \item 62.5 \unit{\micro\gram\per\milli\litre}
    \item 31.25 \unit{\micro\gram\per\milli\litre}
    \item 15.625 \unit{\micro\gram\per\milli\litre}
\end{enumerate}
Now the normal measurement for a large number of concentration is each will be triplicate. so 21 times for each and for standard and control 42 times requires.

Items requires simply to count 1000 \unit{\micro\gram\per\milli\litre}:
 \begin{itemize}
     \item Sample : 10mg   ||  Standard: 10mg
     \item Solvent: 10ml
     \item Precedure (Sample/Standard + Solvent concentration)
 \end{itemize}

 \textbf{ 5\% Bovine Serum Albumin as the Control}:
 The measurement will be on\%w/v which is (g/100ml)* amount 
 now the calculation is :

 \begin{center}
       500 \unit{\micro\liter} of 5\% BSA :
  5\% == 5g/ 100ml 
  here, 5 gram of BSA + water or PBS require to fill 100ml = total= 100ml
  500 \unit{\micro\liter} = 0.5ml
  thus the grams needed :
  (5g/100ml)* 0.5ml = 0.025g == 25mg
 \end{center}
\textbf{procedure:}
\begin{enumerate}
    \item   100 \unit{\micro\litre} water or PBS in the test tube
    \item    25mg of BSA
     \item  300 \unit{\micro\liter} Water or  PBS gently mix
     \item then make 500 \unit{\micro\litre} of the final volume 

\end{enumerate}   


Physiological pH to be maintained Phosphate Buffer Solution is required**  \\
According to the Serial dilution method \textbf{Sample and standard Concentrations}
 \begin{enumerate}
     \item 1000 \unit{\micro\gram\per\milli\litre} : First test tube will be taken 10 mg of sample and 10 ml  of the solvent
   
     \item 500 \unit{\micro\gram\per\milli\litre} :  then  500 \unit{\micro\litre}  of the stock solution from the first test tube then add 500 \unit{\micro\litre} solvent 
     
     \item 250 \unit{\micro\gram\per\milli\litre}  :  500 \unit{\micro\litre} from previous tube then add 500 \unit{\micro\litre} solvent 
     
     \item 125 \unit{\micro\gram\per\milli\litre} : 500 \unit{\micro\litre} from previous tube then add 500 \unit{\micro\litre} solvent
     
     \item 62.5 \unit{\micro\gram\per\milli\litre}  : 500 \unit{\micro\litre} from previous tube then add 500 \unit{\micro\litre}  solvent
     
     \item 31.25 \unit{\micro\gram\per\milli\litre} : take 500 \unit{\micro\litre} from previous then add 500 \unit{\micro\litre} solvent.
     
     \item 15.625 \unit{\micro\gram\per\milli\litre} : take 500 \unit{\micro\litre} from previous tube then add 500 \unit{\micro\litre} solvent

     \item \textbf{Blank: 1000 µL solvent (without extract). }

 \end{enumerate}

Each will now will be require to create triple times which is 21 times total and for control and standard total 42 times.

\subsubsection{\textbf{Procedure}}
When the concentrations are prepared now requires to test with a specific procedure 

\begin{enumerate}     
\item  Taking the control (500 \unit{\micro\litre} of 5\% BSA) incubate first 37 \unit{\degree} for 20 min (Let them to be interact) then 57 \unit{\degree} that makes it 100\% denature at pH (6.3-6.5).
\item Now take the .5ml of the sample concentration then + 4.5 ml of BSA then incubate first 37 \unit{\degree} for 20 min (Let them to be interact) then 57 \unit{\degree} that makes it 100\% denature at pH (6.3-6.5).
    \item A cuvette will be taken where it is must be cleaned.
    \item Spectrophotometer to set 660 nm by clicking in set to \unit{\lambda}.
    \item Cuvette will have 3 ml of Blank (Blank solution) to make remove solvent absorbance rate at 660nm make it zero.
    \item cuvette will be filled with the sample + protein and take absorbance at 660 nm. 
    \item triplicate test for standard same as sample    
\end{enumerate}

\section{Misc. + Discussion }

Now There are some confusions on several researchers that when to maintain of the body physiological pH into In Vitro Lab environment how could it be set. 
There is comparative way to resolve the problem.
\begin{enumerate}
    \item When the Bovine  Serum Albumin  preparing the solution instantly measuring the pH by adding of drop wise of Phosphate Buffer and look at pH(6.3-6.5) ( 1X with 200 \unit{\micro\liter} water + 25mg BSA) Directly.
\end{enumerate}

\begin{figure}
    \centering
    \includegraphics[width=1\linewidth]{In Vitro Anti-Inflammatory Activity Assessment Using Albumin.png}
    \caption{In Vitro Anti-Inflammatory Activity Assessment Using Albumin}
    \label{fig:placeholder}
\end{figure}

\newpage


 
\section{References}

\begin{enumerate}

\item Mizushima, Y., and Kobayashi, M. (1968). 
Interaction of anti-inflammatory drugs with serum proteins, especially albumin. 
\textit{Journal of Pharmacy and Pharmacology}, 20(3), 169--173.

\item Sakat, S., Juvekar, A. R., and Gambhire, M. N. (2010). 
In vitro antioxidant and anti-inflammatory activity of methanol extract of \textit{Oxalis corniculata} Linn. 
\textit{International Journal of Pharmacy and Pharmaceutical Sciences}, 2(1), 146--155.

\item Williams, L. A. D., O'Connar, A., Latore, L., Dennis, O., Ringer, S., Whittaker, J. A., Conrad, J., Vogler, B., and Rosner, H. (2008). 
The in vitro anti-denaturation effects induced by natural products and non-steroidal compounds in heat treated bovine serum albumin. 
\textit{Pharmacognosy Magazine}, 4(15), 1--7.

\end{enumerate}
\end{document}
